% Generated draft: requires mathematical and source review.
\documentclass[11pt]{article}
\usepackage[margin=25mm]{geometry}
\usepackage{fontspec}
\setmainfont{Latin Modern Roman}
\usepackage{amsmath,amssymb,mathtools,unicode-math}
\setmathfont{Latin Modern Math}
\usepackage{enumitem,needspace,xurl,hyperref}
\setlength{\parindent}{0pt}
\setlength{\parskip}{5pt}
\setlength{\emergencystretch}{3em}
\newcommand{\N}{\mathbb{N}}
\newcommand{\Z}{\mathbb{Z}}
\newcommand{\Q}{\mathbb{Q}}
\newcommand{\R}{\mathbb{R}}
\newcommand{\C}{\mathbb{C}}
\newcommand{\F}{\mathbb{F}}
\newcommand{\A}{\mathbb{A}}
\newcommand{\PP}{\mathbb{P}}
\newcommand{\E}{\mathbb{E}}
\newcommand{\eps}{\varepsilon}
\newcommand{\dd}{\,\mathrm{d}}
\begin{document}
\section*{UnsolvedMath ID: 20001120}
\textbf{Source problem code:} catalogue entry
\subsection{Definitions and mathematical statement}
Let \(q\) be any prime power and let \(\F_q\) be the field of order \(q\). A full affine line in \(\F_q^3\) is \(a+\F_q v=\{a+tv:t\in\F_q\}\), where \(v\ne0\). Directions are one-dimensional linear subspaces, so scalar multiples of \(v\) specify the same direction. A finite-field Besicovitch or Kakeya set is \(E\subseteq\F_q^3\) such that for every direction some full affine line in that direction is contained in \(E\).

The damaged source question, attributed to T. Tao, is recovered by the aid as seeking a fixed \(\varepsilon>0\) improving the historical Wolff exponent \(5/2\): prove a lower bound on the scale \(q^{5/2+\varepsilon}\) for every such \(E\). The extraction prints ``Wolf'' and corrupts the final exponent. The literal coefficient-one version \(|E|\geq q^{5/2+\varepsilon}\) should be understood with a sufficiently-large-field qualification if obtained by absorbing an absolute constant.

The supplied aid reports Dvir's polynomial-method solution with the full exponent three:
\[
 |E|\geq\binom{q+2}{3}\geq\frac{q^3}{6}.
\]
Consequently every fixed \(0<\varepsilon<1/2\) gives the literal requested inequality for all sufficiently large prime powers \(q\). Exponent three is optimal in scale because \(|E|\leq q^3\). The aid separately describes a sharper result of Bukh and Chao as
\[
 |E|\geq\frac{q^3+q}{4},
\]
and says that this asymptotically matches constructions. This is a supplied literature assessment, not an independent claim of a newly checked sharp result.

The aid also reports a partial-direction interpolation. Let \(S\subseteq\F_q^2\), and suppose
\[
 \forall(u,v)\in S\ \exists a_{u,v}\in\F_q^3:\quad
 a_{u,v}+\F_q(u,v,1)\subseteq E.
\]
For every integer \(m\) with \(0\leq m\leq q-1\), it claims
\[
 |S|>mq\quad\Longrightarrow\quad |E|\geq\binom{m+3}{3}.
\]
Equivalently, for \(S\ne\varnothing\), putting \(r=\lceil|S|/q\rceil\) gives \(|E|\geq\binom{r+2}{3}\). For \(S=\varnothing\) the direction condition is vacuous and only the trivial lower bound zero is intended. The directions \((u,v,1)\) are precisely this normalized non-horizontal chart; no horizontal directions are required by the partial-direction claim. The strict inequality \(|S|>mq\), integer range, and requirement of full lines are material.
\subsection{Short English statement}
The recovered question asks for a power improvement over the historical \(q^{5/2}\) Kakeya bound in three dimensions. The aid reports Dvir's exponent-three solution, a sharper bound attributed to Bukh--Chao, and an additional exact bound when only some normalized directions are represented.
\subsection{Sources}
\begin{itemize}
\item Local registry: permanent UnsolvedMath ID \(20001120\), exact public code \texttt{catalogue entry}; \texttt{aim-combinatorics-notes.json}, \texttt{notes[244]}.
\item AIM workshop ``Recent trends in additive combinatorics'', Problem 3.8, attributed to T. Tao: \url{https://aimath.org/WWN/additivecomb/additivecomb.pdf}. The supplied verification names \emph{Problems Presented at the Workshop on Recent Trend in Additive Combinatorics}, collected by Ernie Croot and Vsevolod F. Lev, page 8 (PDF index 7), and reports a repaired text layer. The registry nevertheless marks the statement \texttt{reconstructed\_unverified}.
\item The original statement attributes the exponent \(5/2\) to ``Wolf''; the aid identifies the intended name as Wolff. The literature field is empty.
\item Research attempt 1 reports \texttt{full\_solution}, status at run \texttt{solved}, rigor label \texttt{proved}, novelty confidence low. It names Dvir's polynomial method, a sharper result of Bukh and Chao described there as current, and a candidate partial-direction theorem. No full bibliographic entries for Wolff, Dvir or Bukh--Chao are supplied. The current registry status is solved and classification \texttt{SOLVED-IN-LITERATURE}.
\end{itemize}
\end{document}
